\scanpage{109}
\mspage{24}
% archive sheet number: 108
% Lines 1--12: continuation of the large diagonally cancelled block on scan 108.
\deleted{composé de $\deleted{\unsure{\varphi|\varphi_i}}\inserted{\varphi_i}$ avec l'isom précédent est l'inclusion canonique de $R^1f_{i*}(\Lambda_{U_i})$ dans la somme au dernier \unsure{mb|membre} de (***). Donc on en conclut \deleted{\unsure{que}} \deleted{\unsure{\(\varphi\)}} aussitôt que \deleted{\illegible{1}} $E_2^{01}$ est égal au dernier membre de (***), et que $\unsure{\varphi|\varphi_i}$ est un isomorphisme. Cela prouve (i). \deleted{Pour} Supposons maintenant que $X$ satisfait au th. de pureté, et prouvons (ii). On \deleted{\unsure{peut supposer}} \inserted{\unsure{se ramène à}} vérifier l'assertion \deleted{\unsure{fibre par fibre}} \unsure{géométrique}, donc on est ramené \inserted{au cas} $X$ strictement local, et \deleted{\unsure{se trouve à pouvoir}} prouver alors que}

\[
\deleted{\bigwedge^p H^1(U,\Lambda)\longrightarrow H^p(U,\Lambda)}
\]

\deleted{est un isomorphisme.}

% Lines 13--16: live replacement begins below the cancellation boundary.
en d'autres termes que l'hom. canonique

\[
\coprod_{1\leq i\leq r}(\Lambda)_{D_i}\longrightarrow R^1f_*(\deleted{\unsure{\Lambda_X|\Lambda_U}}\unsure{\mu_X|\mu_{nX}}^{\deleted{\unsure{-1|\otimes(-1)}}})
\tag*{(\unsure{4.5.4|4.5.1})}
\]

% Lines 17--23.
est un isomorphisme. On \deleted{\unsure{peut se}} est ramené pour ceci, \deleted{\unsure{aux propriétés}} \inserted{ainsi que pour (ii),} à vérifier l'assertion fibre par fibre, ce qui nous ramène au cas où $X$ est str. local, et à justifier alors que l'hom. \unsure{correspondant à ces}

\[
\deleted{\unsure{\coprod_{1\leq i\leq r}|\coprod_{i\leq r}}\unsure{\Lambda|\mu_n}}\Lambda^r\longrightarrow H^1(U,\mu_n\deleted{\unsure{\Lambda}})
\]

est un isom, ainsi que l'hom.

\[
\bigwedge^d H^1(U,\deleted{\unsure{\mu_n}}\Lambda)\longrightarrow H^d(U,\Lambda)\deleted{\unsure{*|\Lambda}}
\]

dans le cas où $X$ satisfait au th. de pureté cohomologique.
